\documentclass[12pt,a4paper]{article}

% Pacotes essenciais
\usepackage[utf8]{inputenc}
\usepackage[T1]{fontenc}
\usepackage[brazil]{babel}
\usepackage{geometry}
\usepackage{setspace}
\usepackage{mathptmx} 
\usepackage{graphicx}
\usepackage{enumitem}
\usepackage{titlesec}

% Margens ABNT
\geometry{
  a4paper,
  left=3cm,
  right=2cm,
  top=3cm,
  bottom=2cm
}

\begin{document}

\thispagestyle{empty} 

\begin{center}

\textbf{FUNDAÇÃO EDUCACIONAL INACIANA PE. SABÓIA DE MEDEIROS}

\vspace{3cm}

\textbf{LUCAS KERR DO AMARAL} \\
\textbf{MARCELA NALESSO}

\vfill

\textbf{GAME DOCUMENT DESIGN}

\vspace{0.5cm}

A Batalha de Cleitin 2: O Retorno do Homi

\vfill

São Bernardo do Campo, SP \\
2026

\end{center}

\newpage
%%%%%%%%%%%%%%%%%%%%%%%%%%%%%%
% PÁGINA 1 – TÍTULO
%%%%%%%%%%%%%%%%%%%%%%%%%%%%%%

\section*{Itens Iniciais}

\textbf{Sistemas previstos:} PC / Mobile / Console \\[0.3cm]
\textbf{Faixa etária:} 12+ \\[0.3cm]
\textbf{Classificação indicativa (ESRB):} E \\[0.3cm]
\textbf{Data prevista de lançamento:} 2026 \\[2cm]


\newpage

%%%%%%%%%%%%%%%%%%%%%%%%%%%%%%
% PÁGINA 2 – ESBOÇO DO JOGO
%%%%%%%%%%%%%%%%%%%%%%%%%%%%%%

\section*{Esboço do Jogo}

\subsection*{Resumo da História}
Descreva aqui o começo, meio e fim do jogo.

\subsection*{Personagem Principal}
Quem é? Qual sua motivação?

\subsection*{Ângulo de Câmera}
Ex: Terceira pessoa, top-down, 2D lateral.

\subsection*{Fluxo do Jogo}
Explique como o jogador progride.

\subsection*{Condição de Vitória}
Como o jogador vence?

\newpage

%%%%%%%%%%%%%%%%%%%%%%%%%%%%%%
% PÁGINA 3 – PERSONAGENS
%%%%%%%%%%%%%%%%%%%%%%%%%%%%%%

\section*{Personagens}

\subsection*{Protagonista}
História pregressa, personalidade, habilidades.

\subsection*{Habilidades Especiais}
\begin{itemize}
    \item Ataque principal
    \item Habilidade especial
    \item Movimento único
\end{itemize}

\subsection*{Mapa de Controles}
Descreva os comandos principais.

\newpage

%%%%%%%%%%%%%%%%%%%%%%%%%%%%%%
% PÁGINA 4 – GAMEPLAY
%%%%%%%%%%%%%%%%%%%%%%%%%%%%%%

\section*{Gameplay}

\begin{itemize}
    \item Gênero do jogo
    \item Estrutura (fases, capítulos, mundo aberto)
    \item Recursos específicos da plataforma
    \item Modo multiplayer (se houver)
\end{itemize}

\newpage

%%%%%%%%%%%%%%%%%%%%%%%%%%%%%%
% PÁGINA 5 – MUNDO DO JOGO
%%%%%%%%%%%%%%%%%%%%%%%%%%%%%%

\section*{Mundo do Jogo}

\subsection*{Ambientes}
\begin{itemize}
    \item Cidade inicial
    \item Floresta
    \item Fortaleza final
\end{itemize}

\subsection*{Atmosfera}
Descreva música, clima e ambientação.

\newpage

%%%%%%%%%%%%%%%%%%%%%%%%%%%%%%
% PÁGINA 6 – EXPERIÊNCIA DE JOGO
%%%%%%%%%%%%%%%%%%%%%%%%%%%%%%

\section*{Experiência de Jogo}

\begin{itemize}
    \item Sensação principal
    \item Emoções evocadas
    \item Interface e navegação
    \item Design de som e trilha sonora
\end{itemize}

\newpage

%%%%%%%%%%%%%%%%%%%%%%%%%%%%%%
% PÁGINA 7 – MECÂNICAS
%%%%%%%%%%%%%%%%%%%%%%%%%%%%%%

\section*{Mecânicas de Jogo}

\subsection*{Mecânicas Principais}
\begin{itemize}
    \item Plataformas móveis
    \item Power-ups
    \item Sistema de progressão
\end{itemize}

\subsection*{Perigos}
\begin{itemize}
    \item Armadilhas
    \item Obstáculos
\end{itemize}

\newpage

%%%%%%%%%%%%%%%%%%%%%%%%%%%%%%
% PÁGINA 8 – INIMIGOS
%%%%%%%%%%%%%%%%%%%%%%%%%%%%%%

\section*{Inimigos}

\subsection*{Inimigos Comuns}
Características e comportamento.

\subsection*{Chefes (Boss)}
\begin{itemize}
    \item Nome
    \item Ambiente
    \item Estratégia para derrotar
    \item Recompensa
\end{itemize}

\newpage

%%%%%%%%%%%%%%%%%%%%%%%%%%%%%%
% PÁGINA 9 – CUTSCENES
%%%%%%%%%%%%%%%%%%%%%%%%%%%%%%

\section*{Cenas de Corte}

\begin{itemize}
    \item Introdução
    \item Transição de fases
    \item Final do jogo
\end{itemize}

Descreva como serão produzidas (animação, CG, etc).

\newpage

%%%%%%%%%%%%%%%%%%%%%%%%%%%%%%
% PÁGINA 10 – MATERIAL BÔNUS
%%%%%%%%%%%%%%%%%%%%%%%%%%%%%%

\section*{Material Bônus}

\begin{itemize}
    \item Skins desbloqueáveis
    \item Novo modo de jogo
    \item Multiplayer
    \item DLC futura
\end{itemize}

\subsection*{Incentivo para Replay}
Explique por que o jogador jogaria novamente.

\end{document}